% 海洋行星（综述）
% license CCBYSA3
% type Wiki

本文根据 CC-BY-SA 协议转载翻译自维基百科\href{https://en.wikipedia.org/wiki/Ocean_world}{相关文章}。

\begin{figure}[ht]
\centering
\includegraphics[width=6cm]{./figures/76246827a4922672.png}
\caption{地球表面主要由海洋覆盖，海洋占据了地球表面的 $75\%$。因此，地球可以被视为一个水世界，但并非一个完全的海洋行星。} \label{fig_Ocwo_1}
\end{figure}
海洋世界（ocean world）、海洋行星（ocean planet）或水世界（water world）是一类行星或天然卫星，其水圈中包含大量以海洋形式存在的水，这些海洋可以位于表面之下，形成地下海洋，或直接存在于表面，甚至可能淹没所有陆地$^{[1][2][3][4]}$。“海洋世界”一词有时也被用于指代其海洋由其他流体或造海物质（thalassogen）构成的天体$^{[5]}$，例如熔岩（如木卫一 Io）、氨（与水形成共晶混合物，可能存在于土卫六 Titan 的内部海洋中），或碳氢化合物（如土卫六表面的情况，其可能是系外海洋中最常见的一类）$^{[6]}$。对地外海洋的研究被称为行星海洋学（planetary oceanography）。

地球是目前唯一已知在其表面存在液态水体的天文天体，尽管在木星的卫星欧罗巴（Europa）与盖尼米德（Ganymede），以及土星的卫星恩克拉多斯（Enceladus）与土卫六（Titan）上，人们推测其表面之下可能存在地下海洋$^{[7]}$。已经发现若干系外行星具备支持液态水存在的适宜条件$^{[8]}$。此外，在地球上也发现了大量地下水，主要以含水层的形式存在$^{[9]}$。对于系外行星而言，受限于当前技术，尚无法直接观测其表面的液态水，因此通常以大气中的水汽作为替代性指标$^{[10]}$。海洋世界的性质为理解其自身历史以及整个太阳系的形成与演化提供了重要线索；与此同时，它们在生命起源与维持方面的潜力也引起了广泛关注。

2020 年 6 月，NASA 科学家基于数学建模研究报告称，在银河系中拥有海洋的系外行星很可能是普遍存在的$^{[11][12]}$。
\subsection{概述}  
\subsubsection{定义}  
根据 Lunine 的定义，“海洋”被界定为“稳定存在、环绕整个天体的液态水体”$^{[13]}$。此外，“海洋世界”这一术语用于指代太阳系中那些拥有稳定、覆盖全球的液态水体的天体；与之相对，“海洋行星”和“水世界”则专指系外行星（即绕其他恒星运行的行星），其整体组成中含有较高质量分数的水$^{[13]}$。
\subsubsection{太阳系行星天体}
\begin{figure}[ht]
\centering
\includegraphics[width=8cm]{./figures/4b85d59dce5a977e.png}
\caption{土卫二（Enceladus）内部结构示意图} \label{fig_Ocwo_2}
\end{figure}
海洋世界因其可能孕育生命并维持生物活动的潜力而受到天体生物学家的高度关注$^{[4][3]}$。太阳系中被认为拥有地下海洋的主要卫星和矮行星尤为重要，因为与远在数光年之外、受当前技术所限难以直接探测的系外行星不同，这些天体可以通过空间探测器抵达并开展实地研究。除地球之外，太阳系中证据最为充分的水世界包括卡利斯托（Callisto）、恩克拉多斯（Enceladus）、欧罗巴（Europa）、盖尼米德（Ganymede）以及土卫六（Titan）$^{[3][14]}$。其中，欧罗巴和恩克拉多斯由于其外壳较薄且存在冰火山活动，被视为极具吸引力的探测目标。

太阳系中的其他天体也被视为可能拥有地下海洋的候选者，这一判断基于单一类型的观测证据或理论建模结果，包括天卫一（Ariel）$^{[14]}$、天卫三（Titania）$^{[15][16]}$、天卫二（Umbriel）$^{[17]}$、谷神星（Ceres）$^{[3]}$、土卫四（Dione）$^{[18]}$、土卫一（Mimas）$^{[19][20]}$、天卫五（Miranda）$^{[14]}$、天卫四（Oberon）$^{[4][21]}$、冥王星（Pluto）$^{[22]}$、海卫一（Triton）$^{[23]}$、厄里斯（Eris）$^{[4][24]}$以及鸟神星（Makemake）$^{[24]}$。
\subsubsection{系外行星}
\begin{figure}[ht]
\centering
\includegraphics[width=8cm]{./figures/d53b32c9e8bb7ae8.png}
\caption{} \label{fig_Ocwo_3}
\end{figure}
\begin{figure}[ht]
\centering
\includegraphics[width=8cm]{./figures/c7e845b1bfbd19f5.png}
\caption{} \label{fig_Ocwo_4}
\end{figure}